\documentclass[b5paper,11pt]{article}
\usepackage{notestyle}
\usepackage{amsmath}
\usepackage{paracol}

% 避免在未显示编号时出现 subsection 书签锚点重复
\renewcommand{\theHsubsection}{\theHsection.\arabic{subsection}}

% 文档专属：页眉文字
\fancyhead[L]{\fontsize{8pt}{10pt}\selectfont\color{gray!70}心理学通论}

% 添加右下角页脚页数标识
\fancyfoot[R]{\fontsize{8pt}{10pt}\selectfont\color{gray!70} \thepage}

% 文档信息
\title{心理学通论}
\date{2026/04/09}
\author{ClaudeRainer}

\begin{document}
\maketitle

\section{心理学的研究领域和学科性质}

\subsection*{一、心理学的研究领域}
\begin{itemize}
  \item [-] 心理是人脑的\code{机能}，是人脑对\code{客观现实}的\code{主观反映}。
  \item [-] 心理过程是指\code{认知过程}，\code{情感过程}，\code{意志过程}。
\end{itemize}

\begingroup
\setstretch{1.05}
\section{心理学的研究方法}
\setlength{\columnsep}{1em}
\columnratio{0.52,0.48}
\begin{paracol}{2}
  \subsection*{一、一般方法论和研究原则}
  指导心理学研究的哲学方法论就是辩证唯物论和历史唯物论。

  心理学研究具有六大原则：
  \begin{enumerate}
    \item 客观性原则
    \item 系统性原则
    \item 理论与实践相结合
    \item 伦理学原则
    \item 发展性原则
    \item 教育性原则
  \end{enumerate}

  \switchcolumn
  \subsection*{二、具体研究方法}
  \begin{enumerate}
    \item 观察法
    \item 实验法
    \item 调查法
    \item 个案研究法
    \item 内省法
  \end{enumerate}
\end{paracol}
\endgroup

\section{心理学的发展历史与流派}
\hd{一、心理学的发展历史} 1879年，德国心理学家冯特在莱比锡建立了第一个心理实验室，这通常被视为科学心理学诞生的标志。

\hd{心理学的流派}
\begin{itemize}
  \item [-] 第一势力：精神分析学派（弗洛伊德创立，重视潜意识与人格结构）
  \item [-] 第二势力：行为主义学派（华生在巴甫洛夫条件反射学说影响下创立，强调刺激与反应）
  \item [-] 第三势力：人本主义学派（代表人物有罗杰斯、马斯洛，强调自我实现）
  \item [-] 目前主流：认知学派
\end{itemize}

\vspace{1.2em}

\code{心理}：意识、潜意识、前意识

\vspace{0.2em}
\code{人格}：本我、自我、超我

\vspace{0.5em}

\noindent
巴甫洛夫：刺激与信号系统 \\
第一信号系统：具体刺激物；第二信号系统：人类特有的语言和文字

\section{感觉：信息的输入和登录}

\noindent
\textcolor{ctp-green}{阈限分为两种}，分别是绝对阈限和差别阈限。 \\
\textcolor{ctp-green}{感觉剥夺证明} 感觉过程不仅是认识过程的起点，对于维持正常心理功能同样是必不可少的。

\paragraph{知觉的组织}
知觉加工包括对\code{感觉信息}的\code{整合}和赋予加工对象\code{意义}。

\paragraph{自上而下的加工} 借助经验和知识完成对对象的识别，这种符号识别加工受到上下文语义背景以及我们已有知识经验的影响。

\paragraph{自下而上的加工} 把小的知觉单元组合起来，构成更大的单元。

\psplit

\section{注意：信息的选择}

\noindent
\textcolor{ctp-red}{注意} 不是一个独立的心理过程，是一种心理状态。（指向、选择、集中） \\
\textcolor{ctp-red}{注意的功能} 选择性、警觉性、调节性 \\
根据是否需要\textcolor{ctp-red}{意识努力}，注意可以分为\code{有意注意}和\code{无意注意}。

\paragraph{注意的基本品质特点} 注意广度、注意的稳定性、注意的分配性、注意的转移性

\section{记忆：信息的存储和提取}
\paragraph{记忆的分类}

\begin{itemize}
  \item [-] 按内容分类：形象记忆、情绪记忆、运动记忆、语词逻辑记忆
  \item [-] 按保持时间分类：感觉记忆、短时记忆、长时记忆
  \item [-] 按有无自觉努力分类：无意记忆、有意记忆
  \item [-] 按理解在识记中的地位分类：机械记忆、逻辑记忆
  \item [-] 按是否需要努力意志：外显记忆、内隐记忆
\end{itemize}

\begingroup
\setstretch{1.05}
\setlength{\columnsep}{1em}
\columnratio{0.52,0.48}
\begin{paracol}{2}

  \paragraph{遗忘的三种解释}
  \begin{itemize}
    \item [-] 消退说：记忆痕迹随时间消退
    \item [-] 干扰说：新旧信息相互干扰导致遗忘
    \item [-] 抑制说：有意或无意的抑制导致遗忘
  \end{itemize}

  \switchcolumn

  \paragraph{遗忘规律}
  \begin{itemize}
    \item [-] 前摄抑制：旧信息干扰新信息的记忆
    \item [-] 后摄抑制：新信息干扰旧信息的记忆
    \item [-] 位置效应：收尾容易记住，中间容易忘
  \end{itemize}

\end{paracol}
\endgroup

\begingroup
\setstretch{1.05}
\setlength{\columnsep}{1em}
\columnratio{0.52,0.48}
\begin{paracol}{2}
  \subsection*{思维：信息的演算}
  \paragraph{影响问题解决的因素}
  \begin{itemize}
    \item 问题呈现的方式、问题的结构
    \item 定势的影响、功能固着、动机的作用
  \end{itemize}

  \switchcolumn

  \subsection*{心理特征概述}
  \paragraph{个性倾向性}
  \begin{itemize}
    \item 兴趣、需要、动机
    \item 理想、信念、价值观、世界观
  \end{itemize}

\end{paracol}
\endgroup

\paragraph{影响个性心理发展的因素}
\begin{itemize}
  \item 生物因素、身体因素、家庭因素、自然环境因素、学校因素、社会因素、社会文化因素
\end{itemize}

\paragraph{心理特征在教育中的运用}
\begin{itemize}
  \item 有利于\code{因材施教}，促进学生个性发展、培养学生的创新精神和实践能力
  \item 运用于\code{职业指导}，帮助学生了解自己的兴趣、需要、动机，选择适合自己的职业道路
\end{itemize}

\section{个性理论与测验}
\paragraph{荣格的理论} 人格包括三个层次：意识、个人潜意识、集体潜意识

\paragraph{埃里克森的个性发展阶段理论}
\begin{center}
  \begin{tabular}{@{}lllll@{}}
    \toprule
    阶段 & 危机         & 年龄         & 积极品质    & 消极品质    \\ \midrule
    1  & 基本信任与基本不信任 & 0$\sim$1   & 希望      & 恐惧      \\
    2  & 自主对羞怯和疑惑   & 2$\sim$3   & 自我控制和意志 & 自我疑虑    \\
    3  & 主动对内疚      & 4$\sim$5   & 方向和目的   & 无价值感    \\
    4  & 勤奋对自卑      & 6$\sim$11  & 能力      & 无能      \\
    5  & 同一性对角色混乱   & 12$\sim$20 & 忠诚      & 不确定感    \\
    6  & 亲密对孤独      & 20$\sim$24 & 爱       & 两性关系混乱  \\
    7  & 繁殖对停滞      & 25$\sim$64 & 关心      & 自私      \\
    8  & 自我整合对失望    & 65$\sim$死亡 & 明智      & 失望和无意义感 \\ \bottomrule
  \end{tabular}
\end{center}

\begingroup
\setstretch{1.05}
\setlength{\columnsep}{1em}
\columnratio{0.52,0.48}
\begin{paracol}{2}

  \paragraph{马斯洛需求模型（人本主义理论）}
  \begin{itemize}
    \item 自我实现的需要、审美需要
    \item 认知需要、尊重需要
    \item 归属与爱的需要、安全需要
    \item 生理需要
  \end{itemize}

  \switchcolumn

  \paragraph{荣格类型论} 人格类型分为两大类：外向型和内向型。
  \paragraph{霍兰德人格类型} 现实型、研究型、艺术型、社会型、企业型、常规型。
  \paragraph{主要投射测验} 罗夏墨迹测验、主题统觉测验、儿童统觉测验。
\end{paracol}
\endgroup

\paragraph{心理测验的有关指标与使用测验的基本原则}
\begin{itemize}
  \item \textcolor{ctp-red}{信度}：测验结果的稳定性和一致性
  \item \textcolor{ctp-red}{效度}：测验是否测量了它所声称要测量的内容
  \item \textcolor{ctp-red}{常模}：心理测量中的比较标准，在心理测量中常用的标准化样本的分数
  \item 使用测验的基本原则：保护个人隐私、做好准备工作、主试做好充分准备
\end{itemize}

\section{性格特征}
性格是指表现在个体对现实的态度和行为方式中的、比较稳定而独特的心理特征的总和。性格具有以下特点：对现实态度的个性特征、行为方式方面的个性特征，以及明显的社会制约性。

\paragraph{性格的结构特征} 主要包括态度特征、意志特征、情绪特征和理智特征。 \\
其中意志特征主要包括自觉性、坚定性、果断性、自制力和勇敢。

\section{气质}
\paragraph{气质类型}
\begin{itemize}
  \item 多血质：活泼、乐观、外向、热情。情绪不稳定，粗枝大叶。
  \item 黏液质：安静、沉稳、内向、谨慎。死板冷静，缺乏生气。
  \item 胆汁质：易怒、冲动、外向、热情。精力旺盛，表里如一。
  \item 抑郁质：忧郁、悲观、内向、敏感。缓慢单调，不爱交往。
\end{itemize}

\section{心理发展概述}
\begingroup
\setstretch{1.05}
\setlength{\columnsep}{1em}
\columnratio{0.52,0.48}
\begin{paracol}{2}

  \paragraph{儿童心理发展的特征}
  \begin{itemize}
    \item 连续性与阶段性
    \item 不平衡性
    \item 差异性
  \end{itemize}

  \switchcolumn

  \paragraph{心理发展特有的研究方法}
  \begin{itemize}
    \item 横向比较研究
    \item 纵向跟踪研究
    \item 跨文化比较研究
  \end{itemize}

\end{paracol}
\endgroup

\section{认知发展}
\begingroup
\setstretch{1.05}
\setlength{\columnsep}{1em}
\columnratio{0.52,0.48}
\begin{paracol}{2}
  \paragraph{认知发展的阶段}
  \begin{itemize}
    \item 感知运动阶段（出生至2岁）
    \item 前运算阶段（2岁至7岁）
    \item 具体运算阶段（8岁至12岁）
    \item 形式运算阶段（13岁及以上）
  \end{itemize}

  \switchcolumn

  \paragraph{认知发展与平衡化}
  影响认知发展的经典因素主要有三个：成熟、经验和社会环境。

  \paragraph{心理测量流派的智力发展理论（新智力理论）}
  主要包括内部条件（成分）、外部条件（情境）以及连接内外世界的经验。
\end{paracol}
\endgroup

\section{情绪、自我意识和社会性发展}
\paragraph{依恋}
一般认为，6个月至3岁左右是依恋关系发展的重要阶段。依恋类型主要分为三种：
\begin{itemize}
  \item 安全依恋
  \item 焦虑-回避型依恋
  \item 焦虑-矛盾型依恋
\end{itemize}

\paragraph{自我意识的发展} 包括自我认识、自我评价、自我监督和自我调节。

\section{教学设计的理论}
\paragraph{布卢姆教学目标分类} 将教学目标划分为认知、情感和动作技能三个领域。

\section{有效的教学方法}
\noindent
\textcolor{ctp-peach}{先行组织者} \code{奥苏贝尔}：先行组织者是指在学习新材料之前，向学习者提供一个概括性、抽象性、普遍性的框架或概念，以帮助其理解和组织新信息。 \\
\textcolor{ctp-peach}{发现学习与掌握学习} \code{布鲁纳}：发现学习强调通过探索和发现来获得知识，突出学习者的主动参与和创造性思维；掌握学习强调通过系统教学和反复练习来掌握知识与技能，突出教师指导与学习者训练。

\section{认知教学的基本原理}
\noindent
\textcolor{ctp-peach}{皮亚杰认知发展理论} 认知发展阶段分为感知运动阶段、前运算阶段、具体运算阶段和形式运算阶段。\\
\textcolor{ctp-peach}{品德态度} 的形成经历顺从、认同、内化三个阶段。

\psplit

\section{学习类型}
\noindent
\textcolor{ctp-peach}{加涅} 将学习结果分为语言信息、智力技能、认知策略等类型。

\section{学习理论}
\subsection*{学习的联结理论}

\vspace{0.8em}

\noindent
\paragraph{桑代克：联结-试误说}
桑代克把学习看作是通过尝试与错误建立联结的过程。
\begin{itemize}
  \item \code{效果律}：行为的结果会影响该行为在未来发生的可能性。愉快的结果会增强行为，不愉快的结果会削弱行为。
  \item \code{练习律}：反复练习会增强行为之间的联结。
  \item \code{准备律}：学习者只有在准备状态下，学习才更容易发生。
\end{itemize}

\vspace{1.5em}
\noindent
\paragraph{斯金纳：操作性条件作用说}
斯金纳认为，学习是通过操作性条件作用发生的，即个体通过对环境的操作来获得奖励或避免惩罚，
从而增强或削弱特定行为的发生概率。

\subsection*{学习的认知理论}
格式塔学派的学习观是认知学习理论的重要来源，代表人物为\textcolor{ctp-peach}{苛勒}，其理论基础来自对黑猩猩进行的\code{“顿悟实验”}。 \\
认知理论家\textcolor{ctp-peach}{托尔曼}在此基础上提出了“中介变量”的概念；现代认知研究则多以信息加工模型为代表。\textcolor{ctp-peach}{加涅}常被视为将\textcolor{ctp-red}{行为主义学习论}与\textcolor{ctp-red}{认知主义学习论}相结合的代表人物。\\
\textcolor{ctp-red}{整体大于部分之和}。

\subsection*{构建主义学习理论}
构建主义学习理论更加强调学习者的主体作用，强调学习的主动性、社会性和情境性。

在构建主义的学习论中，\textcolor{ctp-teal}{学习是学习者构建自己的知识体系的过程}，
这意味着学习者不是\textcolor{ctp-peach}{被动地接受信息}，
而是通过与环境的互动来\textcolor{ctp-red}{主动构建知识}。构建主义强调学习的社会性，认为学习是在社会互动中发生的，
学习者通过与他人交流和合作来构建知识。
此外，构建主义还强调学习的情境性，认为学习是与特定情境相关的，学习者在特定情境中构建知识。

\section{非智力因素与学习}
\noindent
非智力因素通常指影响学生学习的五种基本心理因素，即\textcolor{ctp-red}{动机、兴趣、情感、意志和性格}。
\paragraph{动机与学习}
\textcolor{ctp-subtext1}{动机是行为的内在动力，是人们在完成任务中力求获得成功的内部动因。} \\
\textcolor{ctp-subtext1}{学习动机是直接推动学生进行学习以达到某种目的的心理动因，即推动学习的主观动力。}

\paragraph{兴趣发展三阶段} \textcolor{ctp-peach}{有趣——乐趣——志趣}

\paragraph{动机中冲突的四种形式}
\begin{itemize}
  \item 双趋冲突：两种积极的选择
  \item 双避冲突：两种消极的选择
  \item 趋避冲突：同一选择具有积极和消极的方面
  \item 多重冲突：多于两种选择，每种选择都有积极和消极的方面
\end{itemize}

\psplit

\section{社会心理}
\begingroup
\setstretch{1.05}
\setlength{\columnsep}{1em}
\columnratio{0.32,0.68}
\begin{paracol}{2}
  \paragraph{影响社会化的因素}
  \begin{itemize}
    \item 家庭
    \item 学校
    \item 伙伴群体
    \item 工作单位
    \item 大众传媒
  \end{itemize}

  \switchcolumn

  \paragraph{印象的定义} 印象是认知者以过去经验为基础，在对认知对象属性进行分析与判断后形成的总体性主观理解。
  \paragraph{印象的特点} 间接性、稳固性、综合性。
  \paragraph{印象形成的模式} 累加模式、平均模式、加权平均模式

\end{paracol}
\endgroup

\section{社会认知的偏差}
\begin{itemize}
  \item \textcolor{ctp-peach}{首因效应和近因效应}：在印象形成过程中，最初接触的信息（首因）和最近接触的信息（近因）对印象的形成有较大的影响。
  \item \textcolor{ctp-peach}{晕轮效应}：对一个人的某个显著特征的印象会影响对其其他特征的评价，导致整体印象被某个特征所主导。
  \item \textcolor{ctp-peach}{社会刻板印象}：对某个群体的成员持有的过于简单化和固定化的看法，可能导致对个体的误解和偏见。
\end{itemize}

\section{社会认知的归因}
个体根据外在的信息和线索来对自己或他人的内在状态或行为原因进行解释和推测的过程，叫做归因。

\textcolor{ctp-peach}{海德的朴素心理学} 将归因分为情境归因和个人倾向归因。
\begin{cquote}
  社会条件、社会舆论、工作的难度、运气的好坏等等因素都属于情境归因的范畴，而能力、努力、性格等因素则属于个人倾向归因的范畴。
\end{cquote}

\textcolor{ctp-peach}{维纳两维度成败归因理论} 维纳将成败归因分为四类：内在稳定、内在不稳定、外在稳定、外在不稳定。
\begin{cquote}
  稳定的内在原因如能力；不稳定的内在原因如努力；稳定的外在原因如任务难度；不稳定的外在原因如运气。
\end{cquote}

\textcolor{ctp-peach}{凯利的归因理论}
\begin{cquote}
  个体通常会依据他人行为的特殊性、一贯性和一致性等信息进行归因。其中，特殊性是指某一具体行为相较于其他相关行为是否具有特殊性。
\end{cquote}

\paragraph{归因的偏差} 基本归因错误、行动者-观察者偏差、防御性归因偏差。

\section{社会影响与群众心理}
主要包括群体心理、从众与服从、去个性化和社会惰化。

\section{与心理健康有关的概念}
\begin{cquote}
  心理辅导是指通过与个体进行交流和指导，帮助其解决心理问题、提升心理素质的一种专业服务。 \\
  心理咨询则是指通过与个体进行深入的交流和分析，帮助其理解和解决心理问题的一种专业服务。\\
  心理治疗则是指通过专业的心理治疗方法，帮助个体解决心理问题、改善心理健康状态的一种专业服务。
\end{cquote}

\textcolor{ctp-red}{出现严重症状、风险行为或超出辅导范围时，应及时转介精神科或临床心理专业人员。}

\end{document}
